% For pdfLaTeX encoding
\usepackage[utf8]{inputenc}
\usepackage[T1]{fontenc}

\usepackage[french]{babel} 
\usepackage{csquotes}
\usepackage{a4wide}
\usepackage{graphicx}
\usepackage{placeins}
\usepackage{tikz}
\usetikzlibrary{chains}
\usetikzlibrary{chains,arrows.meta,positioning,decorations.pathreplacing}
% Custom arrow tip for open circle
\tikzset{
  open arrow/.style = {-{Circle[open]}, shorten <=1pt, shorten >=1pt}
}

\usepackage{pgfplots}
%pour la mise en page des tableaux
\usepackage{array}
\usepackage{tabularx}
\usepackage[table,xcdraw]{xcolor}
\usepackage{wrapfig}
\usepackage{subcaption}
\usepackage{color, colortbl}
\definecolor{Gray}{gray}{0.9}
\definecolor{LightCyan}{rgb}{0.88,1,1}
\usepackage{longtable}
\usepackage{float}
\usepackage{xcolor}
\usepackage{array} % for extrarowheight
\usepackage{lscape}
\usepackage{afterpage}
\usepackage{rotating}
\usepackage[nottoc,notlof,notlot]{tocbibind}
\usepackage{tocloft}
\usepackage{pdfpages}
\graphicspath{../images/}
\newlength\figureheight
\newlength\figurewidth
\usepackage{ifthen}
\usepackage{ifpdf}
\ifpdf
\usepackage[pdftex]{hyperref}
\else
\usepackage{hyperref}
\fi
\usepackage{color}
\hypersetup{%
colorlinks=true,
linkcolor=black,
citecolor=black,
urlcolor=blue}
% \urlstyle{same}
\usepackage[top=2.5cm,bottom=2.5cm,right=2.5cm,left=2.5cm]{geometry}
\usepackage{changepage}

\usepackage{tabularx}
    \newcolumntype{L}{>{\raggedright\arraybackslash}X}
\usepackage{longtable}

\usepackage{xltabular}
\renewcommand{\tablename}{Tableau}
\renewcommand{\figurename}{Figure}

\raggedbottom

%pour les références bibliographiques
\usepackage[style=numeric,sorting=nyt,backend=biber]{biblatex}
\addbibresource{bibliography/references.bib}



% \DefineBibliographyStrings{french}{
%   urlseen = {consulté en}
% }
% \DeclareFieldFormat{urldate}{\mkdate#1}

% \AtEveryBibitem{
%   \ifentrytype{misc}{
%     \clearlist{author}
%     \clearfield{title}
%     \clearfield{note}
%     \clearfield{year}
%     \clearfield{month}
%     \clearfield{day}
%     \clearfield{howpublished}
%     \renewbibmacro*{begentry}{} % supprime l'année avant le lien
%     \renewbibmacro*{finentry}{\par}
%     \renewbibmacro*{url+urldate}{%
%       \printfield{url}%
%       \iffieldundef{urldate}
%         {}
%         { (consulté en \printurldate)}
%     }
%   }{
%     \clearfield{url}
%     \clearfield{urlyear}
%   }
% }

% Only once
\usepackage{xurl}
\usepackage[colorlinks=true,linkcolor=black,citecolor=black,urlcolor=blue,breaklinks]{hyperref}


% pour les algorithmes
\usepackage[linesnumbered,ruled,vlined]{algorithm2e}
\SetKwInput{KwInput}{Input}
\DontPrintSemicolon





\setlength\arrayrulewidth{1pt}
\renewcommand{\baselinestretch}{1.05}
\usepackage{fancyhdr}
\pagestyle{fancy}
\fancyhf{}
\lhead{\bfseries\nouppercase{\leftmark}}
\rfoot{\bfseries\thepage}
\setlength{\headheight}{14.5pt}

\let\headruleORIG\headrule

\renewcommand{\headrule}{\color{black} \headruleORIG}
\renewcommand{\headrulewidth}{1.0pt}
\usepackage{colortbl}
\arrayrulecolor{black}

\fancypagestyle{plain}{
  \fancyhead{}
  \fancyfoot[R]{\bfseries\thepage}
  \renewcommand{\headrulewidth}{0pt}
}



\makeatletter
\def\@textbottom{\vskip \z@ \@plus 1pt}
\let\@texttop\relax
\makeatother

\makeatletter
\def\cleardoublepage{\clearpage\if@twoside \ifodd\c@page\else%
  \hbox{}%
  \thispagestyle{empty}%
  \newpage%
  \if@twocolumn\hbox{}\newpage\fi\fi\fi}
\makeatother

\usepackage{amsthm}
\usepackage{array}
\usepackage{bm}
\usepackage{multirow}
\usepackage[footnote]{acronym}

\usepackage{amssymb,amsmath,bbm,mathtools}
\DeclareMathOperator*{\argmin}{arg\,min}

\usepackage[bottom]{footmisc}


% — permettre la numérotation jusqu’au 4ᵉ niveau (paragraph)
\setcounter{secnumdepth}{4}
% — dans la table des matières, ne pas dépasser le 2ᵉ niveau (subsection)
\setcounter{tocdepth}{2}

% — on va “recycler” \paragraph comme sous-sous-sous-section numérotée
\usepackage{titlesec}

% définir le format du numéro : section.subsection.subsubsection.paragraph
\renewcommand\theparagraph{%
  \thesection.\thesubsection.\thesubsubsection.\arabic{paragraph}%
}

% styler \paragraph pour qu’il ressemble à une vraie section
\titleformat{\paragraph}[hang]
  {\normalfont\normalsize\bfseries}   % police
  {\theparagraph}                     % label
  {1em}                               % espace label–titre
  {}                                  % code avant le titre

% ajuster les espacements verticaux de \paragraph
\titlespacing*{\paragraph}
  {0pt}                               % retrait gauche
  {3.25ex plus 1ex minus .2ex}        % avant
  {1em}                               % après

\newtheoremstyle{break}
  {11pt}{11pt}%
  {\itshape}{}%
  {\bfseries}{}%
  {\newline}{}%
\theoremstyle{break}

%\theoremstyle{definition}
\newtheorem{definition}{Définition}[chapter]

%\theoremstyle{definition}
\newtheorem{theoreme}{Théorème}[chapter]

%\theoremstyle{remark}
\newtheorem{remarque}{Remarque}[chapter]

%\theoremstyle{plain}
\newtheorem{propriete}{Propriété}[chapter]
\newtheorem{exemple}{Exemple}[chapter]

%table break line

\usepackage{makecell}

\renewcommand\theadalign{bl}
\renewcommand\theadfont{\bfseries}
\renewcommand\cellalign{bl}
\renewcommand\cellgape{\Gape[2pt]}

\parskip=5pt
%\sloppy
 %%%%********************************************************************
\definecolor{quotemark}{gray}{0.7}
\makeatletter
\def\fquote{%
    \@ifnextchar[{\fquote@i}{\fquote@i[]}%]
           }%
\def\fquote@i[#1]{%
    \def\tempa{#1}%
    \@ifnextchar[{\fquote@ii}{\fquote@ii[]}%]
                 }%
\def\fquote@ii[#1]{%
    \def\tempb{#1}%
    \@ifnextchar[{\fquote@iii}{\fquote@iii[]}%]
                      }%
\def\fquote@iii[#1]{%
    \def\tempc{#1}%
    \vspace{1em}%
    \noindent%
    \begin{list}{}{%
         \setlength{\leftmargin}{0.1\textwidth}%
         \setlength{\rightmargin}{0.1\textwidth}%
                  }%
         \item[]%
         \begin{picture}(0,0)%
         \put(-15,-5){\makebox(0,0){\scalebox{3}{\textcolor{quotemark}{``}}}}%
         \end{picture}%
         \begingroup\itshape}%
 %%%%********************************************************************
 \def\endfquote{%
 \endgroup\par%
 \makebox[0pt][l]{%
 \hspace{0.8\textwidth}%
 \begin{picture}(0,0)(0,0)%
 \put(15,15){\makebox(0,0){%
 \scalebox{3}{\color{quotemark}''}}}%
 \end{picture}}%
 \ifx\tempa\empty%
 \else%
    \ifx\tempc\empty%
       \hfill\rule{100pt}{0.5pt}\\\mbox{}\hfill\tempa,\ \emph{\tempb}%
   \else%
       \hfill\rule{100pt}{0.5pt}\\\mbox{}\hfill\tempa,\ \emph{\tempb},\ \tempc%
   \fi\fi\par%
   \vspace{0.5em}%
 \end{list}%
 }%
 \makeatother
 %%%%********************************************************************
 
 
 \usepackage{afterpage}

\newcommand\blankpage{%
    \null
    \thispagestyle{empty}%
    \addtocounter{page}{-1}%
    \newpage}
    
\renewcommand{\listalgorithmcfname}{Liste des algorithmes}

\newcommand{\mychapter}[2]{
    
    \chapter*{#2}
    \addcontentsline{toc}{chapter}{#2}
}

\usepackage[page,toc,titletoc,title]{appendix}

\addto\captionsfrench{%
  \renewcommand\appendixname{Annexe}
  \renewcommand\appendixpagename{Annexes}
  \renewcommand{\appendixtocname}{Annexes}
}
\usepackage{etoolbox}
\appto\appendix{\addtocontents{toc}{\protect\setcounter{tocdepth}{0}}}

% reinstate the correct level for list of tables and figures and algorithms
\appto\listoffigures{\addtocontents{lof}{\protect\setcounter{tocdepth}{1}}}
\appto\listoftables{\addtocontents{lot}{\protect\setcounter{tocdepth}{1}}}
\appto\listofalgorithms{\addtocontents{loa}{\protect\setcounter{tocdepth}{1}}}



\usepackage{graphicx}     % pour \includegraphics
\usepackage{booktabs}     % \toprule, \midrule, \bottomrule
\usepackage{subcaption}   % environnement subfigure
\usepackage{siunitx}      % \num, \SI, \percent
\sisetup{detect-all}
\usepackage{pifont}       % symboles ✓ et ✗
\newcommand{\cmark}{\ding{51}}   % ✓
\newcommand{\xmark}{\ding{55}}   % ✗
\DeclareUnicodeCharacter{2009}{\,}



\raggedbottom